%% ============================================================
%% preamble.tex — ARES 현장 진행 보고 (1~4일차)
%% XeLaTeX 컴파일:  xelatex → biber → xelatex → xelatex
%% ============================================================

%% ----- 한국어 지원 -----
\usepackage{kotex}
\usepackage{fontspec}
\defaultfontfeatures{Mapping=tex-text}

%% ----- 폰트 (KoPub 계열 — 격식화: 본문 바탕체(serif) + 제목/UI 돋움체(sans)) -----
%% 격식 있는 조판을 위해 본문(main)을 바탕체(KoPubBatang, serif)로 두고,
%% 제목·머리글·칩 등 \sffamily 로 지정된 UI 요소는 돋움체(KoPubDotum, sans)로 대비를 준다.
%% 굵기는 Medium(보통)/Bold(굵게), 이탤릭 페이스가 없어 \emph/\textit 는 AutoFakeSlant 로 기울인다.
%% 설치: https://www.kopus.org (KoPubWorld/KoPub 글꼴). macOS 도 동일 이름.
\setmainfont{KoPubBatang Medium}[BoldFont={KoPubBatang Bold}, AutoFakeSlant=0.2]
\setsansfont{KoPubDotum Medium}[BoldFont={KoPubDotum Bold}, AutoFakeSlant=0.2]
%% 모노(코드)용 — Consolas(Windows/Office)가 있으면 사용, 없으면 Menlo(macOS 기본)로
%% 폴백해 크로스플랫폼 빌드를 보장한다. (Consolas 부재 시 코드 글자가 누락되던 문제 방지)
\IfFontExistsTF{Consolas}{\setmonofont{Consolas}[Scale=0.86]}{\setmonofont{Menlo}[Scale=0.80]}
\setmainhangulfont{KoPubBatang Medium}[BoldFont={KoPubBatang Bold}, AutoFakeSlant=0.2, Scale=1.0]
\setsanshangulfont{KoPubDotum Medium}[BoldFont={KoPubDotum Bold}, AutoFakeSlant=0.2]
\setmonohangulfont{KoPubDotum Medium}[Scale=0.86]
%% 이모지(흑백) — 부록(코드 분석 전문)이 \emj{}와 raw 이모지를 사용.
%% xetexko가 상위 평면 이모지를 한자(hanja) 부류로 라우팅하므로 한자 폰트를
%% Noto Emoji로 지정해 🚗🔊🚀 등을 렌더한다(본문엔 한자 없음 → 안전).
\newfontfamily\emojifont{NotoEmoji-Regular.ttf}[Path=./fonts/, Scale=0.9]
\setmainhanjafont{NotoEmoji-Regular.ttf}[Path=./fonts/, Scale=0.9]
\setmonohanjafont{NotoEmoji-Regular.ttf}[Path=./fonts/, Scale=0.82]
\newcommand{\emj}[1]{{\emojifont #1}}

%% ----- 레이아웃 -----
\usepackage[a4paper, margin=25mm, headheight=22pt]{geometry}
\usepackage{setspace}
%% 전체 보고서 행간을 기존 1.18 의 1.5배(=1.77)로 확대.
%% (코드 박스 minted 는 아래 etoolbox 훅으로 단일 행간을 유지한다.)
\setstretch{1.77}
\usepackage{microtype}
\emergencystretch=3em
\hfuzz=3pt

%% ----- 색상 -----
\usepackage[table]{xcolor}
\definecolor{aresorange}{HTML}{FA5D29}
\definecolor{deepblue}{HTML}{1A3B5D}
\definecolor{accentcopper}{HTML}{B5651D}
\definecolor{inkdark}{HTML}{222222}
\definecolor{rulegray}{HTML}{B8C0CC}
\definecolor{codebg}{HTML}{F5F6F8}
\definecolor{codekw}{HTML}{0B5394}
\definecolor{codecmt}{HTML}{6A8759}
\definecolor{codestr}{HTML}{A31515}
\definecolor{tablehead}{HTML}{F2E5DE}
\definecolor{chipbg}{HTML}{F2F2F2}

%% ----- 표 -----
\usepackage{booktabs}
\usepackage{longtable}
\usepackage{array}
\usepackage{tabularx}
\usepackage{multirow}
\renewcommand{\arraystretch}{1.25}
\newcolumntype{L}[1]{>{\raggedright\arraybackslash}p{#1}}
\newcolumntype{C}[1]{>{\centering\arraybackslash}p{#1}}
% API 시그니처용(부록 코드 분석): 등폭 작은 글꼴
\newcolumntype{S}[1]{>{\ttfamily\footnotesize\raggedright\arraybackslash}p{#1}}

%% ----- 코드 리스팅 -----
%% ----- 코드 리스팅: minted (Pygments 기반 구문 강조) -----
%% minted 는 fvextra 위에서 조판하므로 XeLaTeX(kotex)에서도 한글 코멘트 옆 괄호가
%% 밀리지 않으며(listings 의 문장부호 재배치 버그 없음) 구문 강조까지 제공한다.
%% 빌드 필수: (1) -shell-escape, (2) Python 패키지 Pygments + latexminted 설치,
%%            (3) 그 실행 파일(pygmentize·latexminted)이 PATH 에 있을 것.
%%            build.sh/build.ps1 참고.
\usepackage[newfloat]{minted}
\usemintedstyle{friendly}
\setminted{
  fontsize=\footnotesize,
  breaklines=true,
  breakanywhere=true,
  tabsize=4,
  frame=single,
  framesep=5pt,
  rulecolor=\color{rulegray},
  numbersep=6pt,
  breaksymbolleft={},
}
%% 코드 리스팅 캡션 라벨을 '코드' 로, 목록 제목을 '코드 목차' 로.
\SetupFloatingEnvironment{listing}{name=코드, listname={코드 목차}}
%% 본문 행간(1.77)을 확대해도 코드 박스 내부 줄간은 기존(단일 행간)을 유지한다.
%% minted 환경 시작 시 \singlespacing 을 걸어, 이후 minted 가 발동하는 폰트 크기
%% 명령(\footnotesize 등)이 늘어난 행간을 물려받지 않게 한다. \inputminted 를 쓰는
%% \srcfile(부록 I) 은 매크로 자체에서 \singlespacing 그룹으로 감싼다.
\usepackage{etoolbox}
\AtBeginEnvironment{minted}{\singlespacing}
%% book 클래스는 그림·표·코드 번호를 장별(예: '부록 A.15')로 매기는데, 부록 장 번호가
%% '부록 A' 라 번호가 지나치게 길어지고 목록에서 겹친다. 문서 전체 연속 번호(그림 1·표 1·
%% 코드 1)로 되돌려 조판을 단순화한다.
\usepackage{chngcntr}
\counterwithout{figure}{chapter}
\counterwithout{table}{chapter}
\AtBeginDocument{\counterwithout{listing}{chapter}}

%% 짧은 인라인 명령/식별자 — 긴 토큰(dknife/ARES2, evaluateValueBlock 등)이
%% 넘칠 때만 문자 사이에서 줄바꿈되도록 seqsplit 으로 감싼다(overfull 방지).
%% 제목→목차(.toc)·PDF 북마크 같은 이동 인자에서도 안전하도록 robust 로 선언한다.
\usepackage{seqsplit}
\DeclareRobustCommand{\code}[1]{{\ttfamily\small\seqsplit{#1}}}

%% 이미지 테두리 박스 — \fbox 여백(\fboxsep·\fboxrule)만큼 이미지 폭을 줄여
%% \linewidth 초과(overfull \hbox 6.8pt)를 방지한다.
\newcommand{\figbox}[1]{\fbox{\includegraphics[width=\dimexpr\linewidth-2\fboxsep-2\fboxrule\relax]{#1}}}

%% 마지막 수단으로 약간의 여유 늘림을 허용해 잔여 overfull \hbox 를 줄인다.
\setlength{\emergencystretch}{3em}

%% ----- 제목 스타일 -----
\usepackage{titlesec}
%% ----- 챕터 표제: 표지 페이지 형태 (수동 \@makechapterhead) -----
%% 번호 있는 장·부록은 홀수 페이지에 '표지'로 시작한다. 표지에는 대형 장 번호와 제목,
%% 하단에 요약 박스가 놓이고, 다음 짝수 페이지를 비운 뒤 홀수에서 본문이 시작된다.
%% ([openright] + 홀수 강제 \cleardoublepage 로 장 시작을 홀수로 맞춘다.)
%% 번호 없는 장(목차·참고문헌 등 \chapter*)은 평범한 표제(\@makeschapterhead)를 쓴다.
%% 표지는 홀수 페이지에서 시작(\chapter 의 \cleardoublepage 가 강제)하되,
%% 표지 다음 쪽을 비우지 않고 본문이 바로 이어지게 한다(\clearpage 만).
\newcommand{\afterchaptercover}{\clearpage}
%% 장 요약(표지 하단 박스) — 각 \chapter 앞에서 \chaptersummary{...} 로 설정
\newcommand{\thischaptersummary}{}
\newcommand{\chaptersummary}[1]{\gdef\thischaptersummary{#1}}
%% 표지 대형 번호 블록 — 본문용(기본). \appendix 뒤 ProgressReport 에서 재정의된다.
\newcommand{\covernumberblock}{%
  {\sffamily\normalsize\color{aresorange}제\,\thechapter\,장}\\[0.1em]
  {\sffamily\bfseries\color{deepblue}\fontsize{108}{108}\selectfont\thechapter}}
\makeatletter
\def\@makechapterhead#1{%
  \thispagestyle{empty}%
  \vspace*{0.12\textheight}%
  \begin{center}\covernumberblock\\[0.85em]
    {\color{rulegray}\rule{0.42\linewidth}{1.1pt}}\\[1.1em]
    {\sffamily\bfseries\Huge\color{deepblue}#1}%
  \end{center}%
  \vfill
  \ifx\thischaptersummary\@empty\else
    {\begin{coversummarybox}\thischaptersummary\end{coversummarybox}}%
    \vspace{0.05\textheight}%
  \fi
  \global\let\thischaptersummary\@empty
  \afterchaptercover}
%% 번호 없는 장(목차·참고문헌) — 평범한 표제
\def\@makeschapterhead#1{%
  \vspace*{16pt}{\raggedright\sffamily\bfseries\LARGE\color{deepblue}#1\par}\nobreak\vskip 18pt}
\makeatother
%% 부(部) 표제 — '제 1 부' 형태(한글)
\renewcommand{\thepart}{제\,\arabic{part}\,부}
\titleformat{\part}[display]
  {\centering\sffamily\bfseries\color{deepblue}}
  {\Huge\thepart}{18pt}{\Huge}
\titlespacing*{\part}{0pt}{80pt}{40pt}
\titleformat{\section}
  {\sffamily\Large\bfseries\color{deepblue}}
  {\thesection}{0.6em}{}
\titleformat{\subsection}
  {\sffamily\large\bfseries\color{deepblue}}
  {\thesubsection}{0.5em}{}
\titleformat{\subsubsection}
  {\sffamily\normalsize\bfseries\color{aresorange}}
  {\thesubsubsection}{0.5em}{}
\titlespacing*{\section}{0pt}{1.4em}{0.7em}

%% ----- 머리글/바닥글 -----
\usepackage{fancyhdr}
\pagestyle{fancy}
\fancyhf{}
\renewcommand{\headrulewidth}{0.4pt}
\renewcommand{\headrule}{\hbox to\headwidth{\color{rulegray}\leaders\hrule height \headrulewidth\hfill}}
\fancyhead[L]{\small\sffamily\color{deepblue}2026 여름 제품개발학기 보고서}
%% 우측(장 제목)이 길면 좌측 문서명과 겹치므로 폭을 제한해 말줄임 처리
\usepackage{truncate}
\fancyhead[R]{\small\sffamily\color{rulegray}\truncate{0.56\textwidth}{\leftmark}}
\fancyfoot[C]{\small\thepage}
%% 머리글에는 장 제목만 표시(번호·라벨 제외)
\renewcommand{\chaptermark}[1]{\markboth{#1}{}}

%% ----- 박스 -----
\usepackage[most]{tcolorbox}
%% 일자 요약 박스
\newtcolorbox{summarybox}{
  breakable, enhanced,
  colback=tablehead!50, colframe=aresorange, boxrule=0pt,
  leftrule=3pt, arc=2pt, left=10pt, right=10pt, top=6pt, bottom=6pt,
}
%% 장 표지 하단 요약 박스
\newtcolorbox{coversummarybox}{
  enhanced, breakable=false,
  colback=tablehead!45, colframe=aresorange, boxrule=0pt, leftrule=3.5pt,
  arc=2pt, left=12pt, right=12pt, top=9pt, bottom=9pt,
  fonttitle=\sffamily\bfseries\color{deepblue}, coltitle=deepblue,
  title={이 장의 요약}, before upper={\small\setstretch{1.3}},
}
%% 주의/메모 박스
\newtcolorbox{notebox}[1][주의]{
  breakable, enhanced,
  colback=codebg, colframe=deepblue, boxrule=0.6pt,
  arc=2pt, left=8pt, right=8pt, top=6pt, bottom=6pt,
  fonttitle=\sffamily\bfseries\color{white},
  coltitle=white, colbacktitle=deepblue,
  title={#1}
}

%% 업무 키워드 칩
\newcommand{\chip}[1]{\tcbox[on line, boxsep=1pt, left=4pt, right=4pt, top=1pt,
  bottom=1pt, colback=chipbg, colframe=chipbg, arc=3pt,
  fontupper=\footnotesize\sffamily]{#1}}

%% ----- 부록(코드 분석 전문)이 쓰는 환경 -----
%% 정보 박스
\newtcolorbox{infobox}[1][]{
  colback=codebg, colframe=deepblue, boxrule=0.6pt,
  arc=2pt, left=8pt, right=8pt, top=6pt, bottom=6pt,
  fonttitle=\sffamily\bfseries\color{white},
  coltitle=white, colbacktitle=deepblue, #1
}
%% 분석 요소 안내 박스
\definecolor{teachbg}{HTML}{FBF4EC}
\newtcolorbox{teachbox}[1]{
  breakable, enhanced,
  colback=teachbg, colframe=accentcopper, boxrule=0.7pt,
  arc=2pt, left=8pt, right=8pt, top=6pt, bottom=6pt,
  fonttitle=\sffamily\bfseries\color{white},
  coltitle=white, colbacktitle=accentcopper,
  title={\hspace{2pt}주요 분석 요소 --- #1}
}
\newcommand{\teachhead}[1]{\smallskip\noindent{\sffamily\bfseries\color{accentcopper}#1}\par\nopagebreak}

%% ----- 기타 -----
\usepackage{enumitem}
\setlist{itemsep=2pt, parsep=1pt, topsep=3pt}
\usepackage{graphicx}
\usepackage{float}
%% ----- 플로트 배치 완화: 그림·표를 띄워(floating) 글이 비는 곳이 없게 -----
%% 한 페이지에 더 많은 플로트를 허용하고, 플로트만 있는 페이지를 줄인다.
\setcounter{topnumber}{4}
\setcounter{bottomnumber}{3}
\setcounter{totalnumber}{6}
\renewcommand{\topfraction}{0.92}
\renewcommand{\bottomfraction}{0.85}
\renewcommand{\textfraction}{0.06}
\renewcommand{\floatpagefraction}{0.72}
\renewcommand{\dbltopfraction}{0.92}
\renewcommand{\dblfloatpagefraction}{0.72}
\usepackage{subcaption}

%% ----- 다이어그램(TikZ) — 부록 코드 분석용 -----
\usepackage{tikz}
\usetikzlibrary{arrows.meta, positioning, fit, calc, shapes.geometric,
  backgrounds, shapes.misc}
\tikzset{
  >={Stealth[length=2.4mm]},
  box/.style={draw=deepblue, fill=codebg, rounded corners=2pt,
    align=center, font=\small\sffamily, inner sep=5pt, minimum height=8mm},
  boxw/.style={box, text width=26mm},
  dom/.style={draw=deepblue, very thick, fill=tablehead, rounded corners=3pt,
    align=center, font=\small\sffamily\bfseries, inner sep=6pt, minimum height=10mm},
  hw/.style={draw=accentcopper, fill=white, rounded corners=2pt,
    align=center, font=\footnotesize\sffamily, inner sep=4pt, minimum height=7mm},
  store/.style={draw=accentcopper, fill=codebg, align=center,
    font=\footnotesize\sffamily, inner sep=4pt, minimum height=7mm,
    cylinder, shape border rotate=90, aspect=0.25},
  flow/.style={->, draw=deepblue, thick},
  back/.style={->, draw=accentcopper, thick, dashed},
  dep/.style={->, draw=deepblue!70, semithick},
  lbl/.style={font=\scriptsize\sffamily\itshape, fill=white, inner sep=1.5pt},
  grp/.style={draw=rulegray, dashed, rounded corners=4pt, inner sep=8pt},
}
\newenvironment{diagram}[2]
  {\def\diagcap{#1}\def\diaglab{#2}\begin{figure}[H]\centering\begin{tikzpicture}}
  {\end{tikzpicture}\caption{\diagcap}\label{\diaglab}\end{figure}}

%% 목차 깊이: 장(챕터)·절(section)까지 표시. 번호는 소절(subsection)까지.
\setcounter{tocdepth}{1}
\setcounter{secnumdepth}{2}
%% 목차의 장 번호 상자 폭 — '부록 A' 라벨이 제목과 겹치지 않도록 넉넉히.
\makeatletter
\renewcommand*\l@chapter[2]{%
  \ifnum \c@tocdepth >\m@ne
    \addpenalty{-\@highpenalty}%
    \vskip 1.0em \@plus\p@
    \setlength\@tempdima{3.4em}%
    \begingroup
      \parindent \z@ \rightskip \@pnumwidth
      \parfillskip -\@pnumwidth
      \leavevmode \bfseries\sffamily\color{deepblue}
      \advance\leftskip\@tempdima \hskip -\leftskip
      #1\nobreak\hfil \nobreak\hb@xt@\@pnumwidth{\hss #2}\par
      \penalty\@highpenalty
    \endgroup
  \fi}
\makeatother
\usepackage[hidelinks]{hyperref}
\hypersetup{
  colorlinks=true,
  linkcolor=deepblue,
  urlcolor=aresorange,
  citecolor=aresorange,
}
%% PDF 북마크(hyperref)에서는 \code/\seqsplit 서식을 제거하고 원문만 사용한다.
\pdfstringdefDisableCommands{\def\code#1{#1}\def\seqsplit#1{#1}}
\usepackage{caption}
\captionsetup{font=small, labelfont={bf,sf}, labelsep=period}
\captionsetup[sub]{font=footnotesize}

%% ----- 라벨 한글화 (Contents→목차, Figure→그림, Table→표, Listing→코드) -----
%% kotex/패키지 기본값을 확실히 덮어쓰도록 \AtBeginDocument 에서 적용한다.
\AtBeginDocument{%
  \renewcommand{\contentsname}{목차}%
  \renewcommand{\figurename}{그림}%
  \renewcommand{\tablename}{표}%
  \renewcommand{\listtablename}{표 목차}%
  \renewcommand{\listfigurename}{그림 목차}%
}

%% ----- 부록 표지(별지) 체계 -----
%% 각 부록은 (1) 홀수 페이지에서 표지로 시작하고, (2) 표지 다음 한 페이지를 비운 뒤,
%% (3) 다시 홀수 페이지에서 본문이 시작되도록 한다. 문서가 oneside 이므로
%% \clearoddpage 로 필요 시 빈 페이지를 삽입해 홀수(recto) 페이지를 강제한다.
\newcommand{\clearoddpage}{%
  \clearpage
  \ifodd\value{page}\else
    \hbox{}\thispagestyle{empty}\newpage
  \fi}
%% 완전 백지 한 장(머리글·쪽번호 없음)
\newcommand{\blankpage}{\clearpage\hbox{}\thispagestyle{empty}\clearpage}
%% oneside 라도 [openright] 의 \cleardoublepage 가 홀수(recto) 페이지를 강제하도록 재정의.
%% (기본 oneside 에서는 \cleardoublepage=\clearpage 라 홀수 강제가 안 됨)
\renewcommand{\cleardoublepage}{\clearpage\ifodd\value{page}\else\hbox{}\thispagestyle{empty}\clearpage\fi}
%% \appendixcover{부록 라벨}{제목}{설명}
%%   표지(홀수) → 백지(짝수) → 본문(홀수, 이어지는 \section 이 조판)
\newcommand{\appendixcover}[3]{%
  \clearoddpage
  \thispagestyle{empty}%
  \begingroup
    \centering
    \vspace*{0.20\textheight}%
    {\sffamily\large\color{aresorange} ARES 제품개발학기 진행 보고 \,\textbullet\, 부록}\\[2.4em]
    {\sffamily\fontsize{44}{48}\selectfont\bfseries\color{deepblue} #1}\\[0.7em]
    {\color{rulegray}\rule{0.46\linewidth}{1.1pt}}\\[1.3em]
    {\fontsize{20}{25}\selectfont\bfseries\color{inkdark} #2}\\[1.8em]
    \begin{minipage}{0.70\linewidth}\centering\normalsize\color{inkdark}\setstretch{1.35} #3\par\end{minipage}%
    \par
    \vfill
    {\sffamily\footnotesize\color{rulegray} (주)코리아사이언스 \,/\, 동명대학교 게임공학과}\\[0.6em]
  \endgroup
  \blankpage
}

%% ----- 참고 문헌 (biblatex + biber) -----
\usepackage[backend=bibtex, style=numeric, sorting=none, giveninits=true]{biblatex}
\addbibresource{references.bib}
